\documentclass[11pt,a4paper]{article}
\usepackage[utf8]{inputenc}
\usepackage[T1]{fontenc}
\usepackage[spanish,es-minimal,es-noshorthands]{babel}
\AtBeginDocument{\renewcommand{\sin}{\operatorname{sen}}\renewcommand{\arcsin}{\operatorname{arcsen}}}
\usepackage{lmodern}
\usepackage[sfdefault,lining]{FiraSans}
\usepackage[a4paper,top=2.3cm,bottom=2.2cm,left=1.8cm,right=1.8cm,headheight=1.3cm,headsep=0.4cm,footskip=1cm]{geometry}
\usepackage{amsmath,amssymb}
\usepackage{icomma}
\usepackage{graphicx}
\usepackage[table]{xcolor}
\usepackage{tikz}
\usepackage{fontawesome5}
\usepackage{ccicons}
\usepackage{fancyhdr}
\usepackage{lastpage}
\usepackage{tasks}
\usepackage{needspace}
\usepackage[most]{tcolorbox}
\usepackage{array}
\usepackage[hidelinks]{hyperref}
\AtBeginDocument{\hypersetup{pdftitle={Ficha de ejercicios de la unidad 8: Límites y continuidad. 1º Bach.},pdfsubject={Matemáticas I 1º Bach.},pdfkeywords={límites, límites en el infinito, comparación de infinitos, indeterminaciones, límites laterales, continuidad, tipos de discontinuidad, funciones a trozos, parámetros, asíntotas, ramas parabólicas}}}

\definecolor{navy}{HTML}{1B2748}
\definecolor{accent}{HTML}{E8603A}
\definecolor{teal}{HTML}{1F9E90}
\definecolor{soft}{HTML}{EEF1F8}
\definecolor{muted}{HTML}{5D6685}

\setlength{\parindent}{0pt}
\setlength{\parskip}{3pt}
\settasks{label=\alph*), label-format=\bfseries\color{navy}, label-width=1.3em, item-indent=1.8em, column-sep=0.6em, after-item-skip=5pt, before-skip=3pt, after-skip=4pt}

% ---------- capçalera i peu ----------
\newcommand{\llicencia}{\href{https://creativecommons.org/licenses/by-nc-sa/4.0/deed.es}{\ccbyncsa}\ \ Licencia Creative Commons Atribución-NoComercial-CompartirIgual 4.0 Internacional (CC BY-NC-SA 4.0)}
\pagestyle{fancy}
\fancyhf{}
\fancyhead[L]{\raisebox{-0.35\height}{\includegraphics[height=0.95cm]{logo-gva.png}}}
\fancyhead[C]{\small\color{muted}Ficha de ejercicios · Unidad 8: Límites y continuidad · 1\textsuperscript{o} Bach.}
\fancyhead[R]{\raisebox{-0.35\height}{\includegraphics[height=0.95cm]{logo-vermellar.png}}}
\fancyfoot[L]{\scriptsize\color{muted}\llicencia}
\fancyfoot[R]{\small\color{muted}\thepage}
\renewcommand{\headrulewidth}{0.4pt}
\renewcommand{\headrule}{\color{muted!40}\hrule height \headrulewidth}
\fancypagestyle{primera}{\fancyhf{}\renewcommand{\headrulewidth}{0pt}
  \fancyfoot[L]{\scriptsize\color{muted}\llicencia}
  \fancyfoot[R]{\small\color{muted}\thepage}}

% ---------- nivells i exercicis ----------
\newcommand{\estrelles}[1]{{\color{accent}\ifcase#1\or\faStar\faStar[regular]\faStar[regular]\or\faStar\faStar\faStar[regular]\or\faStar\faStar\faStar\fi}}
\newcounter{exer}
\newcommand{\ex}[1]{\par\needspace{4\baselineskip}\medskip\refstepcounter{exer}\noindent{\color{navy}\bfseries\theexer.}\ \ }

\newcommand{\seccio}[2]{\par\needspace{16\baselineskip}\bigskip
  \begin{tcolorbox}[colback=navy,colframe=navy,coltext=white,arc=2mm,boxrule=0pt,left=3mm,right=3mm,top=1.2mm,bottom=1.2mm]
  {\large\bfseries #1\quad #2}
  \end{tcolorbox}\vspace{-2pt}}

\newcommand{\exf}[1]{\par\medskip\refstepcounter{exer}\noindent{\color{navy}\bfseries\theexer.}\ \ }
\providecommand{\gt}{>}\providecommand{\lt}{<}
% quadre de recordatori (només a la fitxa d'exercicis)
\newcommand{\recorda}[1]{\par\smallskip\begin{tcolorbox}[colback=accent!7,colframe=accent!55,arc=2mm,boxrule=0.5pt,left=2.5mm,right=2.5mm,top=1mm,bottom=1mm]\small{\color{accent}\faLightbulb[regular]}\ \textbf{Recuerda:}\ #1\end{tcolorbox}}
\newcommand{\res}[1]{#1}\newcommand{\rest}[1]{#1}

\def\fmneed{11}% que el títol d'una secció no quede sol al final del full
% =====================================================================
% Capçalera FentMates de les fitxes: un color per a cada curs.
% Els colors són els de la pàgina inicial (LEVELS a assets/js/materials.js);
% el logo ve de _eines/logo/pdf/ (python3 _eines/logo/genera.py).
%
% Ús, al final del preàmbul de la unitat:
%   % =====================================================================
% Capçalera FentMates de les fitxes: un color per a cada curs.
% Els colors són els de la pàgina inicial (LEVELS a assets/js/materials.js);
% el logo ve de _eines/logo/pdf/ (python3 _eines/logo/genera.py).
%
% Ús, al final del preàmbul de la unitat:
%   % =====================================================================
% Capçalera FentMates de les fitxes: un color per a cada curs.
% Els colors són els de la pàgina inicial (LEVELS a assets/js/materials.js);
% el logo ve de _eines/logo/pdf/ (python3 _eines/logo/genera.py).
%
% Ús, al final del preàmbul de la unitat:
%   \input{../comu/fentmates.tex}
%   \fmcurs{1eso}
% i, al començament del document:
%   \capcalera{Matemàtiques}{1r ESO}{Fitxa d'exercicis · Unitat 2}{Divisibilitat}
% Opcional abans de \input: \def\fmneed{14} (espai mínim abans de cada secció).
% =====================================================================
\makeatletter
\def\fm@logodir{../../logo/pdf}
\providecommand{\fmneed}{10}
\def\fm@cursnom{}

% clau del curs, color i color del text damunt del color
\newcommand{\fm@curs}[3]{\@namedef{fm@color@#1}{#2}\@namedef{fm@text@#1}{#3}}
\fm@curs{1eso}{E6B84A}{navy}
\fm@curs{2eso}{3B82F6}{white}
\fm@curs{3eso}{FF7A50}{white}
\fm@curs{4eso-a}{A78BFA}{white}
\fm@curs{4eso-b}{6D28D9}{white}
\fm@curs{1batx-mi}{9B1C1C}{white}
\fm@curs{1batx-ccss}{EF6B6B}{white}
\fm@curs{2batx-mii}{1E6B47}{white}
\fm@curs{2batx-ccss}{2BB3A3}{white}

\newcommand{\fmcurs}[1]{%
  \@ifundefined{fm@color@#1}{\PackageError{fentmates}{Curs desconegut: #1}{}}{}%
  \edef\fm@tmp{\noexpand\definecolor{curs}{HTML}{\@nameuse{fm@color@#1}}%
               \noexpand\colorlet{curstext}{\@nameuse{fm@text@#1}}}\fm@tmp
  \colorlet{cursfosc}{curs!72!black}%
  \def\fm@logo{\fm@logodir/vertical-#1.pdf}%
  \def\fm@marca{\fm@logodir/marca-#1.pdf}}

% ---------- capçalera de la primera pàgina ----------
% #1 matèria (petit), #2 curs, #3 tipus de document i unitat, #4 títol
\newcommand{\capcalera}[4]{%
  \gdef\fm@cursnom{#2}%
  \par\noindent
  \begin{tcolorbox}[colback=curs,colframe=curs,boxrule=0pt,arc=2.2mm,width=3.5cm,height=2.15cm,
      valign=center,left=3mm,right=2mm,top=1mm,bottom=1mm,nobeforeafter,box align=center]
    \color{curstext}{\footnotesize #1}\par\vspace{1pt}{\fontsize{19}{21}\selectfont\bfseries #2\par}%
  \end{tcolorbox}\hspace{5mm}%
  \begin{minipage}[c]{\dimexpr\linewidth-3.5cm-5mm-2.9cm\relax}
    {\color{cursfosc}\small\bfseries\MakeUppercase{#3}}\par\vspace{2pt}
    {\color{navy}\fontsize{20}{23}\selectfont\bfseries #4\par}
  \end{minipage}\hfill
  \begin{minipage}[c]{2.6cm}\raggedleft\includegraphics[width=2.6cm]{\fm@logo}\end{minipage}\par
  \vspace{2pt}}

% ---------- barres de secció del color del curs ----------
\renewcommand{\seccio}[2]{\par\needspace{\fmneed\baselineskip}\bigskip
  \begin{tcolorbox}[colback=curs,colframe=curs,coltext=curstext,arc=2mm,boxrule=0pt,left=3mm,right=3mm,top=1.2mm,bottom=1.2mm]
  {\large\bfseries #1\quad #2}
  \end{tcolorbox}\vspace{-2pt}}

% ---------- estrelles de dificultat del color del curs (a Batxillerat no se'n posen) ----------
\@ifundefined{estrelles}{}{%
  \renewcommand{\estrelles}[1]{{\color{curs}\ifcase#1\or\faStar\faStar[regular]\faStar[regular]\or\faStar\faStar\faStar[regular]\or\faStar\faStar\faStar\fi}}}

% ---------- capçalera i peu de les pàgines ----------
% primera pàgina: logos institucionals al peu
\fancypagestyle{primera}{\fancyhf{}\renewcommand{\headrulewidth}{0pt}
  \fancyfoot[L]{\raisebox{-0.32\height}{\includegraphics[height=0.78cm]{logo-gva.png}}\hspace{5mm}%
    \raisebox{-0.32\height}{\includegraphics[height=0.78cm]{logo-vermellar.png}}}
  \fancyfoot[R]{\scriptsize\color{muted}\href{https://creativecommons.org/licenses/by-nc-sa/4.0/deed.ca}{\ccbyncsa}\ \ CC BY-NC-SA 4.0 · fentmates.org\qquad\small\thepage}}
% la resta: el símbol a l'esquerra i el curs a la dreta, amb el títol al centre com abans
\fancyhead[L]{\raisebox{-0.3\height}{\includegraphics[height=0.62cm]{\fm@marca}}}
\fancyhead[R]{\small\bfseries\color{cursfosc}\fm@cursnom}
\renewcommand{\headrule}{\color{curs!45}\hrule height \headrulewidth}
\makeatother

%   \fmcurs{1eso}
% i, al començament del document:
%   \capcalera{Matemàtiques}{1r ESO}{Fitxa d'exercicis · Unitat 2}{Divisibilitat}
% Opcional abans de \input: \def\fmneed{14} (espai mínim abans de cada secció).
% =====================================================================
\makeatletter
\def\fm@logodir{../../logo/pdf}
\providecommand{\fmneed}{10}
\def\fm@cursnom{}

% clau del curs, color i color del text damunt del color
\newcommand{\fm@curs}[3]{\@namedef{fm@color@#1}{#2}\@namedef{fm@text@#1}{#3}}
\fm@curs{1eso}{E6B84A}{navy}
\fm@curs{2eso}{3B82F6}{white}
\fm@curs{3eso}{FF7A50}{white}
\fm@curs{4eso-a}{A78BFA}{white}
\fm@curs{4eso-b}{6D28D9}{white}
\fm@curs{1batx-mi}{9B1C1C}{white}
\fm@curs{1batx-ccss}{EF6B6B}{white}
\fm@curs{2batx-mii}{1E6B47}{white}
\fm@curs{2batx-ccss}{2BB3A3}{white}

\newcommand{\fmcurs}[1]{%
  \@ifundefined{fm@color@#1}{\PackageError{fentmates}{Curs desconegut: #1}{}}{}%
  \edef\fm@tmp{\noexpand\definecolor{curs}{HTML}{\@nameuse{fm@color@#1}}%
               \noexpand\colorlet{curstext}{\@nameuse{fm@text@#1}}}\fm@tmp
  \colorlet{cursfosc}{curs!72!black}%
  \def\fm@logo{\fm@logodir/vertical-#1.pdf}%
  \def\fm@marca{\fm@logodir/marca-#1.pdf}}

% ---------- capçalera de la primera pàgina ----------
% #1 matèria (petit), #2 curs, #3 tipus de document i unitat, #4 títol
\newcommand{\capcalera}[4]{%
  \gdef\fm@cursnom{#2}%
  \par\noindent
  \begin{tcolorbox}[colback=curs,colframe=curs,boxrule=0pt,arc=2.2mm,width=3.5cm,height=2.15cm,
      valign=center,left=3mm,right=2mm,top=1mm,bottom=1mm,nobeforeafter,box align=center]
    \color{curstext}{\footnotesize #1}\par\vspace{1pt}{\fontsize{19}{21}\selectfont\bfseries #2\par}%
  \end{tcolorbox}\hspace{5mm}%
  \begin{minipage}[c]{\dimexpr\linewidth-3.5cm-5mm-2.9cm\relax}
    {\color{cursfosc}\small\bfseries\MakeUppercase{#3}}\par\vspace{2pt}
    {\color{navy}\fontsize{20}{23}\selectfont\bfseries #4\par}
  \end{minipage}\hfill
  \begin{minipage}[c]{2.6cm}\raggedleft\includegraphics[width=2.6cm]{\fm@logo}\end{minipage}\par
  \vspace{2pt}}

% ---------- barres de secció del color del curs ----------
\renewcommand{\seccio}[2]{\par\needspace{\fmneed\baselineskip}\bigskip
  \begin{tcolorbox}[colback=curs,colframe=curs,coltext=curstext,arc=2mm,boxrule=0pt,left=3mm,right=3mm,top=1.2mm,bottom=1.2mm]
  {\large\bfseries #1\quad #2}
  \end{tcolorbox}\vspace{-2pt}}

% ---------- estrelles de dificultat del color del curs (a Batxillerat no se'n posen) ----------
\@ifundefined{estrelles}{}{%
  \renewcommand{\estrelles}[1]{{\color{curs}\ifcase#1\or\faStar\faStar[regular]\faStar[regular]\or\faStar\faStar\faStar[regular]\or\faStar\faStar\faStar\fi}}}

% ---------- capçalera i peu de les pàgines ----------
% primera pàgina: logos institucionals al peu
\fancypagestyle{primera}{\fancyhf{}\renewcommand{\headrulewidth}{0pt}
  \fancyfoot[L]{\raisebox{-0.32\height}{\includegraphics[height=0.78cm]{logo-gva.png}}\hspace{5mm}%
    \raisebox{-0.32\height}{\includegraphics[height=0.78cm]{logo-vermellar.png}}}
  \fancyfoot[R]{\scriptsize\color{muted}\href{https://creativecommons.org/licenses/by-nc-sa/4.0/deed.ca}{\ccbyncsa}\ \ CC BY-NC-SA 4.0 · fentmates.org\qquad\small\thepage}}
% la resta: el símbol a l'esquerra i el curs a la dreta, amb el títol al centre com abans
\fancyhead[L]{\raisebox{-0.3\height}{\includegraphics[height=0.62cm]{\fm@marca}}}
\fancyhead[R]{\small\bfseries\color{cursfosc}\fm@cursnom}
\renewcommand{\headrule}{\color{curs!45}\hrule height \headrulewidth}
\makeatother

%   \fmcurs{1eso}
% i, al començament del document:
%   \capcalera{Matemàtiques}{1r ESO}{Fitxa d'exercicis · Unitat 2}{Divisibilitat}
% Opcional abans de \input: \def\fmneed{14} (espai mínim abans de cada secció).
% =====================================================================
\makeatletter
\def\fm@logodir{../../logo/pdf}
\providecommand{\fmneed}{10}
\def\fm@cursnom{}

% clau del curs, color i color del text damunt del color
\newcommand{\fm@curs}[3]{\@namedef{fm@color@#1}{#2}\@namedef{fm@text@#1}{#3}}
\fm@curs{1eso}{E6B84A}{navy}
\fm@curs{2eso}{3B82F6}{white}
\fm@curs{3eso}{FF7A50}{white}
\fm@curs{4eso-a}{A78BFA}{white}
\fm@curs{4eso-b}{6D28D9}{white}
\fm@curs{1batx-mi}{9B1C1C}{white}
\fm@curs{1batx-ccss}{EF6B6B}{white}
\fm@curs{2batx-mii}{1E6B47}{white}
\fm@curs{2batx-ccss}{2BB3A3}{white}

\newcommand{\fmcurs}[1]{%
  \@ifundefined{fm@color@#1}{\PackageError{fentmates}{Curs desconegut: #1}{}}{}%
  \edef\fm@tmp{\noexpand\definecolor{curs}{HTML}{\@nameuse{fm@color@#1}}%
               \noexpand\colorlet{curstext}{\@nameuse{fm@text@#1}}}\fm@tmp
  \colorlet{cursfosc}{curs!72!black}%
  \def\fm@logo{\fm@logodir/vertical-#1.pdf}%
  \def\fm@marca{\fm@logodir/marca-#1.pdf}}

% ---------- capçalera de la primera pàgina ----------
% #1 matèria (petit), #2 curs, #3 tipus de document i unitat, #4 títol
\newcommand{\capcalera}[4]{%
  \gdef\fm@cursnom{#2}%
  \par\noindent
  \begin{tcolorbox}[colback=curs,colframe=curs,boxrule=0pt,arc=2.2mm,width=3.5cm,height=2.15cm,
      valign=center,left=3mm,right=2mm,top=1mm,bottom=1mm,nobeforeafter,box align=center]
    \color{curstext}{\footnotesize #1}\par\vspace{1pt}{\fontsize{19}{21}\selectfont\bfseries #2\par}%
  \end{tcolorbox}\hspace{5mm}%
  \begin{minipage}[c]{\dimexpr\linewidth-3.5cm-5mm-2.9cm\relax}
    {\color{cursfosc}\small\bfseries\MakeUppercase{#3}}\par\vspace{2pt}
    {\color{navy}\fontsize{20}{23}\selectfont\bfseries #4\par}
  \end{minipage}\hfill
  \begin{minipage}[c]{2.6cm}\raggedleft\includegraphics[width=2.6cm]{\fm@logo}\end{minipage}\par
  \vspace{2pt}}

% ---------- barres de secció del color del curs ----------
\renewcommand{\seccio}[2]{\par\needspace{\fmneed\baselineskip}\bigskip
  \begin{tcolorbox}[colback=curs,colframe=curs,coltext=curstext,arc=2mm,boxrule=0pt,left=3mm,right=3mm,top=1.2mm,bottom=1.2mm]
  {\large\bfseries #1\quad #2}
  \end{tcolorbox}\vspace{-2pt}}

% ---------- estrelles de dificultat del color del curs (a Batxillerat no se'n posen) ----------
\@ifundefined{estrelles}{}{%
  \renewcommand{\estrelles}[1]{{\color{curs}\ifcase#1\or\faStar\faStar[regular]\faStar[regular]\or\faStar\faStar\faStar[regular]\or\faStar\faStar\faStar\fi}}}

% ---------- capçalera i peu de les pàgines ----------
% primera pàgina: logos institucionals al peu
\fancypagestyle{primera}{\fancyhf{}\renewcommand{\headrulewidth}{0pt}
  \fancyfoot[L]{\raisebox{-0.32\height}{\includegraphics[height=0.78cm]{logo-gva.png}}\hspace{5mm}%
    \raisebox{-0.32\height}{\includegraphics[height=0.78cm]{logo-vermellar.png}}}
  \fancyfoot[R]{\scriptsize\color{muted}\href{https://creativecommons.org/licenses/by-nc-sa/4.0/deed.ca}{\ccbyncsa}\ \ CC BY-NC-SA 4.0 · fentmates.org\qquad\small\thepage}}
% la resta: el símbol a l'esquerra i el curs a la dreta, amb el títol al centre com abans
\fancyhead[L]{\raisebox{-0.3\height}{\includegraphics[height=0.62cm]{\fm@marca}}}
\fancyhead[R]{\small\bfseries\color{cursfosc}\fm@cursnom}
\renewcommand{\headrule}{\color{curs!45}\hrule height \headrulewidth}
\makeatother

\fmcurs{1batx-mi}

\begin{document}
\thispagestyle{primera}

\capcalera{Matemáticas I}{1\textsuperscript{o} Bach.}{Ficha de ejercicios · Unidad 8}{\raggedright Límites y continuidad}
\vspace{10pt}
Nombre y apellidos: \rule{6.9cm}{0.4pt}\hfill Grupo: \rule{1.6cm}{0.4pt}\hfill Fecha: \rule{2.2cm}{0.4pt}
\vspace{8pt}

\begin{tcolorbox}[colback=soft,colframe=soft,arc=2mm,left=3mm,right=3mm,top=1.5mm,bottom=1.5mm]
\small Resuelve los ejercicios en el cuaderno mostrando todos los pasos: antes de transformar una expresión, sustituye y comprueba si hay de verdad una indeterminación. Las secciones siguen el orden de los apuntes. Puedes comprobar los resultados con la calculadora de límites de los apuntes interactivos.
\end{tcolorbox}

% =====================================================
\seccio{1}{Límites en el infinito}

\exf{1} Observa las gráficas y escribe, para cada una, el límite cuando $x\to+\infty$ y cuando $x\to-\infty$. Las rectas discontinuas son asíntotas.\begin{tasks}(4)
  \task \par\begin{tikzpicture}[scale=0.3,baseline=(current bounding box.center)]\begin{scope}\draw[muted!30,very thin] (-5,-4) grid (5,4);\draw[-latex,thick] (-5,0)--(5.4,0);\draw[-latex,thick] (0,-4)--(0,4.4);\clip (-5,-4) rectangle (5,4);\draw[accent,line width=1.5pt,smooth,samples=90,domain=-3.6:3.6] plot (\x,{\x*\x*\x/8-\x});\end{scope}\end{tikzpicture}
  \task \par\begin{tikzpicture}[scale=0.3,baseline=(current bounding box.center)]\begin{scope}\draw[muted!30,very thin] (-5,-4) grid (5,4);\draw[-latex,thick] (-5,0)--(5.4,0);\draw[-latex,thick] (0,-4)--(0,4.4);\clip (-5,-4) rectangle (5,4);\draw[accent,line width=1.5pt,smooth,samples=90,domain=-5:2.5] plot (\x,{2^\x-2});\draw[teal,dashed,thick] (-5,-2)--(5,-2);\end{scope}\end{tikzpicture}
  \task \par\begin{tikzpicture}[scale=0.3,baseline=(current bounding box.center)]\begin{scope}\draw[muted!30,very thin] (-5,-4) grid (5,4);\draw[-latex,thick] (-5,0)--(5.4,0);\draw[-latex,thick] (0,-4)--(0,4.4);\clip (-5,-4) rectangle (5,4);\draw[accent,line width=1.5pt,smooth,samples=90,domain=-5:5] plot (\x,{2*\x*\x/(\x*\x+1)});\draw[teal,dashed,thick] (-5,2)--(5,2);\end{scope}\end{tikzpicture}
  \task \par\begin{tikzpicture}[scale=0.3,baseline=(current bounding box.center)]\begin{scope}\draw[muted!30,very thin] (-5,-4) grid (5,4);\draw[-latex,thick] (-5,0)--(5.4,0);\draw[-latex,thick] (0,-4)--(0,4.4);\clip (-5,-4) rectangle (5,4);\draw[accent,line width=1.5pt,smooth,samples=90,domain=-5:5] plot (\x,{1.5*sin(2*\x r)});\end{scope}\end{tikzpicture}
\end{tasks}\recorda{Para $x\to+\infty$ mira el \textbf{final derecho} de la gráfica y para $x\to-\infty$, el \textbf{final izquierdo}. Si la gráfica oscila sin acercarse a ningún valor, el límite no existe.}

\ex{2} Calcula, para cada función, $\displaystyle\lim_{x\to-\infty}f(x)$ y $\displaystyle\lim_{x\to+\infty}f(x)$. Fíjate bien en el dominio de las raíces.\begin{tasks}(2)
  \task $f(x)=-2x^3+7x^2$
  \task $f(x)=3x^4-7x$
  \task $f(x)=10^x$
  \task $f(x)=\sqrt{5x-8}$
  \task $f(x)=\sqrt{-2x^2+1}$
  \task $f(x)=-5^x$
  \task $f(x)=\sqrt{3-x}$
  \task $f(x)=\dfrac{x^2+3}{-x^3}$
  \task $f(x)=\dfrac{-x^3}{x^2+3}$
  \task $f(x)=\dfrac{5x^3-10}{3x^3+10x^2}$
\end{tasks}\recorda{Potencias, exponenciales de base mayor que 1, raíces y logaritmos tienden a $+\infty$ cuando $x\to+\infty$. Un polinomio se comporta como su término de mayor grado. Si la función no está definida hacia un infinito, el límite no existe.}

\ex{1} Calcula los límites siguientes comparando los grados del numerador y del denominador:\begin{tasks}(3)
  \task $\displaystyle\lim_{x\to +\infty}\dfrac{x^2-5x+3}{3x-5}$
  \task $\displaystyle\lim_{x\to +\infty}\dfrac{x^2+3}{x^3}$
  \task $\displaystyle\lim_{x\to +\infty}\dfrac{3x^2-5x+1}{2x^2-6}$
  \task $\displaystyle\lim_{x\to -\infty}\dfrac{4x^3-x}{2-x^2}$
  \task $\displaystyle\lim_{x\to -\infty}\dfrac{1-6x}{3x+2}$
  \task $\displaystyle\lim_{x\to -\infty}\dfrac{2x^2-1}{5-x}$
\end{tasks}\recorda{$\displaystyle\lim_{x\to\pm\infty}\frac{P(x)}{Q(x)}=\lim_{x\to\pm\infty}\frac{a_nx^n}{b_mx^m}$: \ si $n>m$, $\pm\infty$ (hay que mirar el signo); \ si $n=m$, $\frac{a_n}{b_m}$; \ si $n<m$, $0$.}

\ex{2} Calcula los límites siguientes con exponenciales y raíces:\begin{tasks}(3)
  \task $\displaystyle\lim_{x\to +\infty}\left(\tfrac13\right)^x$
  \task $\displaystyle\lim_{x\to -\infty}\left(3^x+1\right)$
  \task $\displaystyle\lim_{x\to -\infty}0{,}4^x$
  \task $\displaystyle\lim_{x\to +\infty}\dfrac{\sqrt{4x^2+1}}{x+3}$
  \task $\displaystyle\lim_{x\to -\infty}\dfrac{\sqrt{x^2+1}}{2x}$
  \task $\displaystyle\lim_{x\to +\infty}\left(e^{-x}+2\right)$
\end{tasks}


% =====================================================
\seccio{2}{Comparación de infinitos}

\exf{1} Ordena estos infinitos de menor a mayor orden cuando $x\to+\infty$: \quad $x^2,\ \ 3^x,\ \ \ln x,\ \ x^5,\ \ \sqrt x,\ \ 1{,}1^x$.\recorda{Cuando $x\to+\infty$: \ logaritmos $\ll$ potencias $\ll$ exponenciales (de base mayor que 1). Entre potencias gana el exponente mayor y entre exponenciales, la base mayor.}

\ex{2} Calcula los límites siguientes comparando infinitos:\begin{tasks}(3)
  \task $\displaystyle\lim_{x\to +\infty}\dfrac{x^4}{2^x}$
  \task $\displaystyle\lim_{x\to +\infty}\dfrac{3^x}{x^{20}}$
  \task $\displaystyle\lim_{x\to +\infty}\dfrac{\ln x}{x}$
  \task $\displaystyle\lim_{x\to +\infty}\left(x^3-e^x\right)$
  \task $\displaystyle\lim_{x\to +\infty}\dfrac{2^x+x^2}{2^{x+1}}$
  \task $\displaystyle\lim_{x\to +\infty}\left(\sqrt x-\ln x\right)$
  \task $\displaystyle\lim_{x\to +\infty}\left(x-2^x\right)$
  \task $\displaystyle\lim_{x\to +\infty}\dfrac{5^x+x}{4^x}$
\end{tasks}

\ex{2} Indeterminación $\infty-\infty$ por comparación de infinitos:\begin{tasks}(2)
  \task $\displaystyle\lim_{x\to +\infty}\left(x^2-\sqrt{x^5+3}\right)$
  \task $\displaystyle\lim_{x\to +\infty}\left(3^x-\sqrt{x^3-2}\right)$
\end{tasks}


% =====================================================
\seccio{3}{Indeterminaciones en el infinito}

\exf{2} Indeterminación $\frac{\infty}{\infty}$ con raíces:\begin{tasks}(2)
  \task $\displaystyle\lim_{x\to +\infty}\dfrac{\sqrt{x^2+1}+2x}{3x-1}$
  \task $\displaystyle\lim_{x\to +\infty}\dfrac{\sqrt{4x^4+x}-x}{x^2+5}$
  \task $\displaystyle\lim_{x\to +\infty}\dfrac{3x+\sqrt x}{x-\sqrt x}$
  \task $\displaystyle\lim_{x\to +\infty}\dfrac{\sqrt{4x^3+x-2}}{8x^2-7}$
\end{tasks}\recorda{Nos quedamos con el término de mayor grado de cada parte: \ $\sqrt{x^2+1}\approx x$ (si $x>0$), \ $\sqrt{4x^3}=2x^{3/2}$.}

\ex{2} Indeterminación $\infty-\infty$ con fracciones:\begin{tasks}(2)
  \task $\displaystyle\lim_{x\to +\infty}\left(\dfrac{2x^2-5x}{x+3}-2x\right)$
  \task $\displaystyle\lim_{x\to +\infty}\left(\dfrac{x^2}{x-2}-\dfrac{x^2}{x+2}\right)$
  \task $\displaystyle\lim_{x\to +\infty}\left(\dfrac{x^3}{x^2-1}-x\right)$
  \task $\displaystyle\lim_{x\to +\infty}\left(\dfrac{x^2+1}{x}+\dfrac{5x-x^2}{3x-1}\right)$
\end{tasks}\recorda{Si la indeterminación $\infty-\infty$ viene de una suma o resta de fracciones, \textbf{hacemos la operación} con denominador común y calculamos el límite de la fracción que resulta.}

\ex{2} Indeterminación $\infty-\infty$ con raíces:\begin{tasks}(3)
  \task $\displaystyle\lim_{x\to +\infty}\left(x-\sqrt{x^2-1}\right)$
  \task $\displaystyle\lim_{x\to +\infty}\left(\sqrt{x+2}-\sqrt{x-2}\right)$
  \task $\displaystyle\lim_{x\to +\infty}\left(\sqrt{x^2+6x}-x\right)$
  \task $\displaystyle\lim_{x\to +\infty}\left(\sqrt{x^2-2}-\sqrt{x^2+x}\right)$
  \task $\displaystyle\lim_{x\to +\infty}\left(2x-\sqrt{4x^2-3x}\right)$
\end{tasks}\recorda{Multiplicamos y dividimos por el \textbf{conjugado}: $(A-B)(A+B)=A^2-B^2$ hace desaparecer las raíces del numerador.}

\ex{2} Indeterminación $1^{\infty}$:\begin{tasks}(3)
  \task $\displaystyle\lim_{x\to +\infty}\left(\dfrac{x}{x+2}\right)^{2x+1}$
  \task $\displaystyle\lim_{x\to +\infty}\left(\dfrac{2x-2}{2x+1}\right)^{3x^2}$
  \task $\displaystyle\lim_{x\to +\infty}\left(1+\dfrac3x\right)^{2x}$
  \task $\displaystyle\lim_{x\to +\infty}\left(\dfrac{x+1}{x-2}\right)^{x}$
  \task $\displaystyle\lim_{x\to +\infty}\left(\dfrac{x^2+1}{x^2-3}\right)^{x^2}$
  \task $\displaystyle\lim_{x\to +\infty}\left(\dfrac{3x+1}{x+2}\right)^{x}$
\end{tasks}\recorda{Si $\lim f(x)=1$ y $\lim g(x)=\pm\infty$: \ $\lim f(x)^{g(x)}=e^{\lim\,(f(x)-1)\cdot g(x)}$. Si la base no tiende a 1, no hay indeterminación.}

\ex{2} Problemas. Calcula el límite que responde a cada pregunta:\begin{tasks}(1)
  \task El coste medio (en euros) de fabricar $x$ unidades de un producto es $C(x)=\dfrac{12x+3\,000}{x}$. ¿A qué valor se acerca el coste por unidad si se fabrican muchísimas?
  \task La concentración de un medicamento en sangre $t$ horas después de tomarlo es $c(t)=\dfrac{8t}{t^2+4}$ (mg/l). ¿Qué pasa con la concentración a muy largo plazo?
  \task Una población de conejos evoluciona según $P(t)=\dfrac{500t+300}{t+1}$, con $t$ en meses. ¿Cuál es la población inicial? ¿Tiene algún límite el número de conejos?
\end{tasks}


% =====================================================
\seccio{4}{Límites en un punto}

\exf{1} \label{ex:graf}Observa la gráfica de la función $g$ y calcula, si existen, los límites que se indican. Calcula también $g(-2)$, $g(1)$ y $g(3)$.\par\smallskip{\centering \begin{tikzpicture}[scale=0.4,baseline=(current bounding box.center)]\begin{scope}\draw[muted!30,very thin] (-6,-5) grid (6,5);\draw[-latex,thick] (-6,0)--(6.4,0);\draw[-latex,thick] (0,-5)--(0,5.4);\clip (-6,-5) rectangle (6,5);\draw[accent,line width=1.5pt,smooth,samples=90,domain=-6:-2.18] plot (\x,{1+1/(\x+2)});\draw[accent,line width=1.5pt,smooth,samples=90,domain=-2:1] plot (\x,{\x+1});\draw[accent,line width=1.5pt,smooth,samples=90,domain=1:5.9] plot (\x,{(\x-3)*(\x-3)/2+1});\draw[teal,dashed,thick] (-6,1)--(-2,1);\draw[teal,dashed,thick] (-2,-5)--(-2,5);\filldraw[fill=white,draw=accent,line width=1.1pt] (-2,-1) circle (0.17);\fill[accent] (1,2) circle (0.17);\filldraw[fill=white,draw=accent,line width=1.1pt] (1,3) circle (0.17);\filldraw[fill=white,draw=accent,line width=1.1pt] (3,1) circle (0.17);\fill[accent] (3,-1) circle (0.17);\node[below,font=\scriptsize] at (1,0) {$1$};\node[below,font=\scriptsize] at (3,0) {$3$};\node[above left,font=\scriptsize] at (-2,0) {$-2$};\node[left,font=\scriptsize] at (0,1) {$1$};\end{scope}\end{tikzpicture}\par}\smallskip\begin{tasks}(4)
  \task $\displaystyle\lim_{x\to -\infty}g(x)$
  \task $\displaystyle\lim_{x\to +\infty}g(x)$
  \task $\displaystyle\lim_{x\to -2^-}g(x)$
  \task $\displaystyle\lim_{x\to -2^+}g(x)$
  \task $\displaystyle\lim_{x\to 1^-}g(x)$
  \task $\displaystyle\lim_{x\to 1^+}g(x)$
  \task $\displaystyle\lim_{x\to 1}g(x)$
  \task $\displaystyle\lim_{x\to 3}g(x)$
\end{tasks}\recorda{Un punto vacío quiere decir que ese punto no es de la gráfica. El límite en un punto existe si los dos límites laterales coinciden, y no depende del valor de la función en el punto.}

\ex{1} Comprueba, dando a $x$ valores cada vez más próximos al punto, que $\displaystyle\lim_{x\to1^-}\frac{1}{x-1}=-\infty$ y que $\displaystyle\lim_{x\to-1^-}\left(x^2+5\right)=6$.

\ex{1} Calcula los límites siguientes. Cuando al sustituir salga $\frac k0$, estudia los límites laterales:\begin{tasks}(3)
  \task $\displaystyle\lim_{x\to 2}(x^3-2x+1)$
  \task $\displaystyle\lim_{x\to -1}\dfrac{x^2+3}{x-1}$
  \task $\displaystyle\lim_{x\to 3}\dfrac{2x}{x-3}$
  \task $\displaystyle\lim_{x\to 0}\dfrac{x-1}{x^2}$
  \task $\displaystyle\lim_{x\to 1}\dfrac{x+2}{1-x}$
  \task $\displaystyle\lim_{x\to 2}\dfrac{1-x}{(x-2)^2}$
  \task $\displaystyle\lim_{x\to 3^+}\ln(x-3)$
\end{tasks}\recorda{Si al sustituir sale $\frac{k}{0}$ con $k\neq0$, el límite es infinito: para saber su signo, sustituimos un valor un poco más pequeño y otro un poco más grande que el punto.}

\ex{2} Calcula los límites de las funciones a trozos siguientes en los puntos indicados:\begin{tasks}(2)
  \task \[f(x)=\begin{cases}2x+5&\text{si }x<3\\[3pt]-x+2&\text{si }x\geq3\end{cases}\]Calcula $\displaystyle\lim_{x\to1}f(x)$, $\displaystyle\lim_{x\to5}f(x)$ y $\displaystyle\lim_{x\to3}f(x)$.
  \task \[g(x)=\begin{cases}x^2+1&\text{si }x<-1\\[3pt]1-x&\text{si }-1\leq x<2\\[3pt]x-3&\text{si }x\geq2\end{cases}\]Calcula $\displaystyle\lim_{x\to-1}g(x)$ y $\displaystyle\lim_{x\to2}g(x)$.
\end{tasks}


% =====================================================
\seccio{5}{Indeterminaciones en un punto}

\exf{2} Resuelve las indeterminaciones $\frac00$ factorizando:\begin{tasks}(3)
  \task $\displaystyle\lim_{x\to 1}\dfrac{x-1}{x^2-3x+2}$
  \task $\displaystyle\lim_{x\to 3}\dfrac{x^3-9x}{x^3-3x^2}$
  \task $\displaystyle\lim_{x\to 2}\dfrac{x^2-4}{x^2+x-6}$
  \task $\displaystyle\lim_{x\to 1}\dfrac{x^3-1}{x^2-1}$
  \task $\displaystyle\lim_{x\to -2}\dfrac{x^2+2x}{x^2+4x+4}$
  \task $\displaystyle\lim_{x\to 1}\dfrac{x^3-3x+2}{x^3-x^2-x+1}$
\end{tasks}\recorda{Si $P(a)=Q(a)=0$, los dos polinomios son divisibles por $(x-a)$: factorizamos (factor común, identidades notables, ecuación de segundo grado o Ruffini), simplificamos y volvemos a sustituir.}

\ex{2} Indeterminaciones $\frac00$ con raíces:\begin{tasks}(2)
  \task $\displaystyle\lim_{x\to 9}\dfrac{\sqrt x-3}{x-9}$
  \task $\displaystyle\lim_{x\to 4}\dfrac{\sqrt{x+5}-3}{x-4}$
  \task $\displaystyle\lim_{x\to 1}\dfrac{x-1}{\sqrt{x+3}-2}$
  \task $\displaystyle\lim_{x\to 3}\dfrac{2-\sqrt{x+1}}{x-3}$
\end{tasks}\recorda{Multiplicamos y dividimos por el conjugado de la expresión que tiene la raíz y simplificamos el factor que se anula.}

\ex{2} Indeterminación $1^\infty$ en un punto:\begin{tasks}(2)
  \task $\displaystyle\lim_{x\to 0}(1+3x)^{\frac2x}$
  \task $\displaystyle\lim_{x\to 2}(x-1)^{\frac{1}{x-2}}$
  \task $\displaystyle\lim_{x\to 1}\left(\dfrac{x+3}{4}\right)^{\frac{2}{x-1}}$
  \task $\displaystyle\lim_{x\to 1}(2-x)^{\frac{x}{x-1}}$
\end{tasks}\recorda{La fórmula es la misma que en el infinito: \ $\displaystyle\lim_{x\to a}f(x)^{g(x)}=e^{\lim\limits_{x\to a}(f(x)-1)\cdot g(x)}$.}

\ex{2} Calcula $\displaystyle\lim_{x\to1}\left(\frac{1}{x-1}-\frac{2}{x^2-1}\right)$. ¿Qué indeterminación aparece?


% =====================================================
\seccio{6}{Continuidad y tipos de discontinuidad}

\exf{1} Clasifica las discontinuidades de estas gráficas en el punto $x=c$:\begin{tasks}(4)
  \task \par\begin{tikzpicture}[scale=0.32,baseline=(current bounding box.center)]\begin{scope}\draw[muted!30,very thin] (-3,-3) grid (3,3);\draw[-latex,thick] (-3,0)--(3.4,0);\draw[-latex,thick] (0,-3)--(0,3.4);\clip (-3,-3) rectangle (3,3);\draw[accent,line width=1.5pt,smooth,samples=90,domain=-3:0.66] plot (\x,{1/(\x-1)});\draw[accent,line width=1.5pt,smooth,samples=90,domain=1.34:3] plot (\x,{1/(\x-1)});\draw[teal,dashed,thick] (1,-3)--(1,3);\node[below,font=\scriptsize] at (1,0) {$c$};\end{scope}\end{tikzpicture}
  \task \par\begin{tikzpicture}[scale=0.32,baseline=(current bounding box.center)]\begin{scope}\draw[muted!30,very thin] (-3,-3) grid (3,3);\draw[-latex,thick] (-3,0)--(3.4,0);\draw[-latex,thick] (0,-3)--(0,3.4);\clip (-3,-3) rectangle (3,3);\draw[accent,line width=1.5pt,smooth,samples=90,domain=-3:1] plot (\x,{\x+1});\draw[accent,line width=1.5pt,smooth,samples=90,domain=1:3] plot (\x,{\x-2});\fill[accent] (1,2) circle (0.17);\filldraw[fill=white,draw=accent,line width=1.1pt] (1,-1) circle (0.17);\node[below,font=\scriptsize] at (1,0) {$c$};\end{scope}\end{tikzpicture}
  \task \par\begin{tikzpicture}[scale=0.32,baseline=(current bounding box.center)]\begin{scope}\draw[muted!30,very thin] (-3,-3) grid (3,3);\draw[-latex,thick] (-3,0)--(3.4,0);\draw[-latex,thick] (0,-3)--(0,3.4);\clip (-3,-3) rectangle (3,3);\draw[accent,line width=1.5pt,smooth,samples=90,domain=-2.4:2.4] plot (\x,{\x*\x/2-1});\filldraw[fill=white,draw=accent,line width=1.1pt] (1,-0.5) circle (0.17);\fill[accent] (1,1.5) circle (0.17);\node[below,font=\scriptsize] at (1,0) {$c$};\end{scope}\end{tikzpicture}
  \task \par\begin{tikzpicture}[scale=0.32,baseline=(current bounding box.center)]\begin{scope}\draw[muted!30,very thin] (-3,-3) grid (3,3);\draw[-latex,thick] (-3,0)--(3.4,0);\draw[-latex,thick] (0,-3)--(0,3.4);\clip (-3,-3) rectangle (3,3);\draw[accent,line width=1.5pt,smooth,samples=90,domain=-3:3] plot (\x,{\x/2+1});\filldraw[fill=white,draw=accent,line width=1.1pt] (1,1.5) circle (0.17);\node[below,font=\scriptsize] at (1,0) {$c$};\end{scope}\end{tikzpicture}
\end{tasks}\recorda{$f$ es continua en $x=c$ si $\exists f(c)$, $\exists\displaystyle\lim_{x\to c}f(x)$ y ambos coinciden. Si el límite existe pero falla alguna otra condición, la discontinuidad es \textbf{evitable}; si los laterales son números diferentes, de \textbf{salto finito}; si un lateral es infinito y el otro un número, de \textbf{salto infinito}; si los dos son infinitos, \textbf{asintótica}.}

\ex{1} Vuelve a la gráfica de $g$ del ejercicio \ref{ex:graf}. ¿En qué puntos no es continua? Clasifica cada discontinuidad.

\ex{1} Estudia la continuidad de cada función en el punto indicado y, si es discontinua, di de qué tipo es:\begin{tasks}(2)
  \task $f(x)=\dfrac{x^2-x}{x}$ \ en $x=0$
  \task $g(x)=\dfrac{x+1}{x-2}$ \ en $x=2$
  \task $h(x)=\begin{cases}x^2+1&\text{si }x<0\\x+1&\text{si }x\geq0\end{cases}$ \ en $x=0$
  \task $k(x)=\begin{cases}x+2&\text{si }x\leq1\\\dfrac{1}{x-1}&\text{si }x>1\end{cases}$ \ en $x=1$
  \task $p(x)=\begin{cases}\dfrac{x^2-4}{x-2}&\text{si }x\neq2\\[2pt]3&\text{si }x=2\end{cases}$ \ en $x=2$
\end{tasks}

\ex{2} Estudia la continuidad de las funciones siguientes:\begin{tasks}(2)
  \task $f(x)=\dfrac{x^2-1}{x^2+x-2}$
  \task $g(x)=\dfrac{x+3}{x^2+4}$
  \task $h(x)=\dfrac{x}{x^2-4x}$
  \task $k(x)=\sqrt{x-4}$
\end{tasks}\recorda{Las funciones racionales son continuas en todo $\mathbb R$ excepto donde se anula el denominador. En esos puntos, si el límite es un número la discontinuidad es evitable y si es $\frac k0$, asintótica.}

\ex{2} Indica si cada una de las funciones siguientes es continua y dónde presenta discontinuidades:\begin{tasks}(2)
  \task $f(x)=\dfrac{x^2+2x}{x+5}$
  \task $g(x)=\dfrac{x^2+2x}{x}$
  \task \[h(x)=\begin{cases}1&\text{si }x<3\\[3pt]4&\text{si }x\geq3\end{cases}\]
  \task \[k(x)=\begin{cases}1&\text{si }x\leq0\\[3pt]1-x^2&\text{si }0<x<2\\[3pt]-3&\text{si }x\geq2\end{cases}\]
\end{tasks}


% =====================================================
\seccio{7}{Funciones a trozos. Parámetros}

\exf{2} Estudia la continuidad de las funciones siguientes:\begin{tasks}(1)
  \task \[f(x)=\begin{cases}x-3&\text{si }x<2\\[3pt]x^2-5&\text{si }x\geq2\end{cases}\]
  \task \[g(x)=\begin{cases}2x+1&\text{si }x<0\\[3pt]x^2+1&\text{si }0\leq x<2\\[3pt]\dfrac{3}{x-1}&\text{si }x\geq2\end{cases}\]
  \task \[h(x)=\begin{cases}x^2&\text{si }x\leq-1\\[3pt]\dfrac1x&\text{si }-1<x<1\\[3pt]\sqrt x&\text{si }x\geq1\end{cases}\]
\end{tasks}\recorda{Primero, la continuidad de cada trozo en su intervalo (¡cuidado con los denominadores y las raíces!). Después, en cada punto donde se rompe la función, comparamos $f(a)$ y los dos límites laterales.}

\ex{2} Escribe $f(x)=\dfrac{|x^2-4|}{x-2}$ como una función a trozos y estudia su continuidad en $x=2$.

\ex{2} Calcula $m$ para que esta función sea continua en todo $\mathbb R$:\[f(x)=\begin{cases}x^2-3&\text{si }x\leq2\\[3pt]-x^2+m&\text{si }x>2\end{cases}\]\recorda{Si la función depende de un parámetro, \textbf{igualamos los límites laterales} en el punto donde se rompe y resolvemos la ecuación.}

\ex{2} Halla el valor del parámetro para que cada función sea continua en todo $\mathbb R$:\begin{tasks}(1)
  \task \[f(x)=\begin{cases}ax+3&\text{si }x<1\\[3pt]x^2-a&\text{si }x\geq1\end{cases}\]
  \task \[g(x)=\begin{cases}\dfrac{x^2-9}{x-3}&\text{si }x\neq3\\[3pt]k&\text{si }x=3\end{cases}\]
  \task \[h(x)=\begin{cases}2^x+a&\text{si }x\leq0\\[3pt]\dfrac{x+4}{x+2}&\text{si }x>0\end{cases}\]
\end{tasks}

\ex{3} Halla $a$ y $b$ para que la función siguiente sea continua en todo $\mathbb R$:\[f(x)=\begin{cases}x+a&\text{si }x<0\\[3pt]bx+1&\text{si }0\leq x\leq2\\[3pt]x^2-a&\text{si }x>2\end{cases}\]

\ex{2} Un taxi cobra $P(x)=3+1{,}5x$ euros por un viaje de $x\leq10$ km y $P(x)=ax+8$ si $x>10$ km. ¿Cuánto debe valer $a$ para que el precio no dé ningún salto en los 10 km? ¿Cuánto cuesta, entonces, un viaje de 25 km?


% =====================================================
\seccio{8}{Ramas infinitas y asíntotas}

\exf{2} Halla las ramas infinitas de las funciones siguientes. Si hay asíntotas, estudia la posición de la curva respecto de ellas:\begin{tasks}(3)
  \task $f(x)=\dfrac{x^2-1}{x^2+4}$
  \task $g(x)=\dfrac{x}{(x-1)^2}$
  \task $h(x)=\dfrac{2x^2}{x-3}$
\end{tasks}\recorda{AV: donde el límite es infinito ($\frac k0$). \ AH: $\displaystyle\lim_{x\to\pm\infty}f(x)=b$. \ AO en una racional cuando el grado de arriba es una unidad mayor: la asíntota es el cociente de la división. La posición se ve con el signo de $f(x)$ menos la asíntota.}

\ex{2} Calcula todas las asíntotas de las funciones siguientes y estudia los límites laterales en cada asíntota vertical:\begin{tasks}(2)
  \task $f(x)=\dfrac{3x+1}{x-2}$
  \task $f(x)=\dfrac{x^2}{x^2-4}$
  \task $f(x)=\dfrac{x+1}{x^2-1}$
  \task $f(x)=\dfrac{x^2+x+1}{x-1}$
  \task $f(x)=\dfrac{x^3+1}{x^2}$
  \task $f(x)=\dfrac{3\cdot2^x}{2^x+1}$
\end{tasks}

\ex{3} Calcula las asíntotas oblicuas de $f(x)=\sqrt{x^2+4x}$ con los límites $m=\displaystyle\lim\frac{f(x)}{x}$ i $n=\displaystyle\lim\,[f(x)-mx]$. Hazlo cuando $x\to+\infty$ y cuando $x\to-\infty$.

\ex{2} ¿Cuáles de estas funciones tienen ramas parabólicas cuando $x\to+\infty$? Justifícalo calculando $\lim f(x)$ i $m=\lim\frac{f(x)}{x}$.\begin{tasks}(4)
  \task $f(x)=x^3-3x$
  \task $g(x)=\ln x$
  \task $h(x)=\dfrac{x^3}{x-1}$
  \task $k(x)=\dfrac{x^2+1}{x}$
\end{tasks}

\ex{3} Halla $a$, $b$ y $c$ para que la función $f(x)=\dfrac{ax+b}{x+c}$ tenga una asíntota vertical en $x=-2$, una asíntota horizontal en $y=3$ y pase por el punto $(0,1)$.

\vfill
\begin{center}\color{muted}\small\faCheckCircle[regular]\ \ ¡Ánimo! Comprueba tus resultados con la calculadora de límites de los apuntes interactivos.\end{center}

\end{document}
